\documentclass[11pt]{article}

\usepackage[a4paper,margin=1in]{geometry}
\usepackage{amsmath,amssymb,amsthm,mathtools}
\usepackage{enumitem}
\usepackage[hidelinks]{hyperref}

\newtheorem{theorem}{Theorem}[section]
\newtheorem{proposition}[theorem]{Proposition}
\newtheorem{lemma}[theorem]{Lemma}
\newtheorem{corollary}[theorem]{Corollary}
\newtheorem{conjecture}[theorem]{Conjecture}
\newtheorem{program}[theorem]{Proof program}
\theoremstyle{definition}
\newtheorem{definition}[theorem]{Definition}
\newtheorem{remark}[theorem]{Remark}

\newcommand{\whp}{\text{with high probability}}
\newcommand{\E}{\mathbb E}
\newcommand{\Pp}{\mathbb P}
\newcommand{\Nn}{\log n}
\newcommand{\gapscale}{\dfrac{n}{(\log n)^3}}
\newcommand{\smallscale}{\dfrac{n}{(\log n)^4}}
\newcommand{\Afour}{A_4}
\newcommand{\Dfour}{D_4}
\newcommand{\rplus}{r_+}
\newcommand{\rco}{r_4^{\mathrm{co}}}

\title{Erd\H{o}s Problem 625: theorem upgrades and follow-up programs}
\author{Research planning note}
\date{26 July 2026}

\begin{document}
\maketitle

\begin{abstract}
This note isolates the strongest theorem statements currently latent in the
candidate proof of Erd\H{o}s Problem 625 and formulates the next substantive
research targets.  Statements are classified as conditional consequences of
the existing proof architecture or as new targets requiring additional work.
No conditional or target statement is represented as already externally
verified.
\end{abstract}

\section{Notation and the root displacement}

Let
\[
 q=\log 2,\qquad N=\log n,\qquad H_n=\frac{n}{N^3}.
\]
For $G_n\sim G(n,1/2)$, write $\chi(G_n)$ for the chromatic number and
$\zeta(G_n)$ for the cochromatic number.

Let $\delta_n\in[0,1)$ be the independence-number phase parameter.  For the
four-size signed support, define
\[
 \Afour(\delta)=\log 2-\Dfour(\delta),
\]
where $\Dfour$ is the limiting entropy loss caused by restricting the deficit
support to four consecutive sizes.

The root calculation used by the candidate proof has the form
\begin{equation}
 \rplus(n)-\rco(n)
 =\left(\frac{q^2}{4}\Afour(\delta_n)+o(1)\right)H_n.
 \label{eq:root-gap}
\end{equation}
Here $\rplus$ is the ordinary-coloring root and $\rco$ is the signed
four-support root.

\section{Conditional main-paper theorems}

\begin{theorem}[Phase-resolved full-sequence gap]
\label{thm:phase-gap}
Assume the normalized signed four-support second moment and the
rare-seed-to-typical-completion step in the candidate proof.  Then
\[
 \chi(G_n)-\zeta(G_n)
 \ge
 \left(\frac{q^2}{8}\Afour(\delta_n)-o(1)\right)H_n
\]
with high probability along the full sequence of integers $n$.
\end{theorem}

\begin{proof}[Proof blueprint]
Place the signed witness at
\[
 k_{1/2}=\left\lceil\frac{\rco+\rplus}{2}\right\rceil.
\]
Equation~\eqref{eq:root-gap} gives
\[
 \rplus-k_{1/2}
 =\frac12(\rplus-\rco)+O(1)
 =\left(\frac{q^2}{8}\Afour(\delta_n)+o(1)\right)H_n.
\]
The normalized second moment gives a rare signed cocoloring seed.  The
amplifier adds $o(H_n)$ classes, while integer rounding also contributes
$o(H_n)$.  Intersect with the ordinary chromatic lower event.
\end{proof}

\begin{corollary}[Certified uniform coefficient]
\label{cor:uniform}
Assume Theorem~\ref{thm:phase-gap} and the certified entropy inequality
\[
 \Afour(\delta)>\log\!\left(\frac{1000}{639}\right)
 \qquad(0\le\delta\le1).
\]
Then
\[
 \chi(G_n)-\zeta(G_n)
 \ge
 \frac{q^2}{8}\log\!\left(\frac{1000}{639}\right)H_n
\]
with high probability.  The coefficient equals
\[
 0.026896409808379\ldots.
\]
\end{corollary}

\begin{corollary}[Simultaneous comparison with the complement]
\label{cor:complement}
Under the hypotheses of Corollary~\ref{cor:uniform},
\[
 \min\{\chi(G_n),\chi(\overline G_n)\}-\zeta(G_n)
 \ge
 \frac{q^2}{8}\log\!\left(\frac{1000}{639}\right)H_n
\]
with high probability.
\end{corollary}

\begin{proof}
Use $\zeta(\overline G)=\zeta(G)$ and the distributional identity
$\overline G_n\stackrel d=G_n$.  Apply Corollary~\ref{cor:uniform} to $G_n$
and $\overline G_n$ and take a union bound.
\end{proof}

\begin{proposition}[Tunable cochromatic upper tail]
\label{prop:tunable-tail}
Suppose a $k_n$-cocoloring seed has probability at least
$\exp(-\Lambda_n)$.  Then there is an absolute constant $C$ such that, for
every deterministic $r>0$,
\[
 \Pp\!\left(
 \zeta(G_n)>k_n+C\left[
  \frac{\sqrt{n\Lambda_n}+\sqrt{nr}}{\log n}+n^{1/3}+1
 \right]\right)
 \le e^{-r}+o(1).
\]
\end{proposition}

\begin{remark}
The proposition should be stated independently of the one value of $r$ used
in the final proof.  It is a general output of the rare-seed amplifier.
\end{remark}

\section{Exact endpoint normalization and the remaining closure}

Fix a full endpoint table $L$.  Let $P$ denote a selected block-level matching
and let every selected cell carry a literal full physical stub matching.
The finite endpoint formalization proves that the aggregate sum over these
decorated objects is exactly the manuscript reference weight $W(L)$.

\begin{proposition}[Decorated endpoint normalization]
\label{prop:endpoint-normalization}
For every feasible four-type endpoint table $L$,
\[
 \sum_{D\in\mathcal D(L)}w_{\mathrm{atom}}(D)=W(L),
\]
where $\mathcal D(L)$ is the family of block pairings decorated by full-cell
stub matchings and $w_{\mathrm{atom}}$ is the common signed reward divided by
the ambient falling factorial.
\end{proposition}

The proof proceeds through the exact cardinality identity
\[
 |\mathcal D(L)|\,
 \Bigl(\prod_{i,j}\ell_{ij}!\Bigr)
 \Bigl(\prod_{i,j}m_{ij}!^{\ell_{ij}}\Bigr)
 =
 R(L)C(L)S(L),
\]
where $R(L),C(L)$ are the row and column block-selection products and $S(L)$
is the product of full-cell stub selections.  All cancellation in
$\mathbb R_{\ge0}\cup\{\infty\}$ is performed only after positivity and
finiteness of the factorial factors are established.

The remaining closure is an exact reindexing of every attained high physical
skeleton by a block support, a deficit vector, and local partial stub
matchings.

For a cell with endpoint sizes $m$ and $m+d$ and deficit $h$, define
\begin{equation}
 R_{m,d}(h)
 =\frac{\binom mh}{(d+1)(d+2)\cdots(d+h)}
 2^{-hm+h(h+1)/2}.
 \label{eq:local-ratio}
\end{equation}
If $J$ is the actual total multiplicity and $H$ the total deficit, then
\begin{equation}
 \frac{(n)_{J+H}}{(n)_J}=(n-J)_H\le n^H.
 \label{eq:global-denominator}
\end{equation}
Consequently the aggregate decorated-to-full comparison is
\begin{equation}
 \frac{w(P,m-h)}{w_{\mathrm{full}}(P)}
 \le\prod_{e\in P}n^{h_e}R_{m_e,d_e}(h_e).
 \label{eq:aggregate-comparison}
\end{equation}

\begin{program}[Global Section VIII closure]
Prove the following in one theorem chain.
\begin{enumerate}[label=(\roman*)]
\item Construct an equivalence between attained high physical skeletons and
triples consisting of a block matching, an admissible deficit vector, and
local partial stub matchings.
\item Sum the local physical fibres to recover the aggregate weight used in
\eqref{eq:aggregate-comparison}.
\item Use the high condition $2h<m$ to prove
\[
 n^hR_{m,d}(h)
 \le\left(\frac{nm}{2^{\lfloor2m/3\rfloor}}\right)^h.
\]
\item Sum the geometric deficit fibres and identify the full reference sum
with $\sum_LW(L)$ using Proposition~\ref{prop:endpoint-normalization}.
\item Apply endpoint transport to obtain
\[
 \operatorname{BareSkeletonSum}_n
 \le\exp\left\{o\left(\frac n{(\log n)^4}\right)\right\}.
\]
\item Compose with the q-only literal attachment theorem to obtain the
normalized second moment.
\end{enumerate}
\end{program}

\section{Near-root placement and removal of the midpoint loss}

For $0<\theta<1$, define
\[
 k_\theta
 =\left\lceil
 \rco+\theta(\rplus-\rco)
 \right\rceil.
\]
Then
\begin{equation}
 \rplus-k_\theta
 =\left((1-\theta)\frac{q^2}{4}\Afour(\delta_n)+o(1)\right)H_n.
 \label{eq:theta-gap}
\end{equation}

\begin{lemma}[Error-envelope criterion]
\label{lem:error-envelope}
Suppose every proof error requiring a positive buffer above the signed root is
$O(n/N^4)$.  If $\theta_n\to0$ and
\[
 \theta_nN\to\infty,
\]
then the buffer
\[
 \theta_n(\rplus-\rco)
 =\Theta\left(\theta_n\frac n{N^3}\right)
\]
dominates the entire $n/N^4$ error scale.
\end{lemma}

\begin{proof}
Divide the buffer by $n/N^4$.  Equation~\eqref{eq:root-gap} gives a ratio of
order $\theta_nN$, which tends to infinity.
\end{proof}

\begin{conjecture}[Near-root placement theorem]
\label{conj:near-root}
The profile construction and normalized second moment are uniform for
placements satisfying the hypothesis of Lemma~\ref{lem:error-envelope}.  In
particular, for $\theta_n=N^{-1/2}$,
\[
 \chi(G_n)-\zeta(G_n)
 \ge
 \left(\frac{q^2}{4}\Afour(\delta_n)-o(1)\right)H_n
\]
with high probability.
\end{conjecture}

The certified fixed coefficient would be
\[
 \frac{q^2}{4}\log\!\left(\frac{1000}{639}\right)
 =0.053792819616758\ldots.
\]

\begin{program}[Uniformity audit for Conjecture~\ref{conj:near-root}]
Expose the placement dependence in:
\begin{enumerate}[label=(\arabic*)]
\item integer profile feasibility and rounding;
\item finite-to-limiting entropy approximation;
\item partial-diagonal estimates;
\item endpoint and all-deficit estimates;
\item q-only residual attachments;
\item Paley--Zygmund seed probability;
\item the final amplifier.
\end{enumerate}
Every occurrence of a fixed positive midpoint margin must be replaced by an
explicit lower bound of order $\theta_nH_n$.
\end{program}

\section{Sharp four-support phase constant}

Let $F_S(T)$ be the support-$S$ constrained entropy value and
$\lambda_S(T)$ its selected multiplier.  The envelope identity is
\[
 \frac{d}{dT}F_S(T)=-\lambda_S(T).
\]
For full support and $S_4=\{2,3,4,5\}$,
\[
 D_4(T)=F_\infty(T)-F_4(T),
\qquad
 D_4'(T)=\lambda_4(T)-\lambda_\infty(T).
\]
If $T=T_0(\delta)$, then
\begin{equation}
 A_4'(\delta)
 =T_0'(\delta)
 \bigl(\lambda_\infty(T_0(\delta))-\lambda_4(T_0(\delta))\bigr).
 \label{eq:phase-derivative}
\end{equation}

\begin{conjecture}[Endpoint minimum]
\label{conj:endpoint-minimum}
The function $A_4$ is minimized at the terminal phase endpoint.  Numerically,
\[
 \min_{0\le\delta\le1}A_4(\delta)
 =0.520701335491\ldots.
\]
\end{conjecture}

\begin{program}[Interval proof of Conjecture~\ref{conj:endpoint-minimum}]
\begin{enumerate}[label=(\roman*)]
\item Prove monotonicity and a rational enclosure for $T_0(\delta)$.
\item Partition the $T$-interval into finitely many rational subintervals.
\item On each interval, enclose $\lambda_4(T)$ and
$\lambda_\infty(T)$ and certify the sign in
\eqref{eq:phase-derivative}.
\item Evaluate the endpoint entropy difference using directed interval
arithmetic.
\end{enumerate}
\end{program}

If certified, the midpoint and near-root coefficients would be respectively
\[
 0.031271565748\ldots,
 \qquad
 0.062543131497\ldots.
\]

\section{Balanced structure of near-optimal cocolorings}

For a proportion $\rho$ of clique classes, the sign-assignment entropy is
\[
 H(\rho)=-\rho\log\rho-(1-\rho)\log(1-\rho).
\]
Relative to unrestricted signs, the loss is
\[
 \log2-H(\rho).
\]
Near $\rho=1/2+x$,
\[
 \log2-H(1/2+x)=2x^2+O(x^4).
\]

\begin{proposition}[Balanced rare seed]
Assume the normalized signed second moment.  Restrict the witness to exactly
$\lfloor k/2\rfloor$ clique labels.  Then the seed probability deteriorates by
at most a polynomial factor in $k$.
\end{proposition}

\begin{proof}
Use
\[
 \binom{k}{\lfloor k/2\rfloor}\ge\frac{2^k}{k+1}
\]
and the pointwise domination of the balanced count by the unrestricted signed
count.  The normalized second moment increases by at most $(k+1)^2$.
\end{proof}

\begin{conjecture}[Balance stability]
\label{conj:balance}
For every fixed $\varepsilon>0$, there exists $c(\varepsilon)>0$ such that,
with high probability, every cocoloring whose clique-class proportion lies
outside $[1/2-\varepsilon,1/2+\varepsilon]$ uses at least
\[
 \rco(n)+c(\varepsilon)H_n
\]
nonempty classes.  Consequently every cocoloring using
$\rco(n)+o(H_n)$ classes has clique proportion $1/2+o(1)$.
\end{conjecture}

\begin{program}[First-moment proof of Conjecture~\ref{conj:balance}]
Refine the signed first moment by the number of clique labels.  Replace the
factor $2^k$ by $\binom{k}{\rho k}$ and repeat the root-displacement analysis
with $\log2$ replaced by $H(\rho)$.  The entropy deficit shifts the root by
order $H_n$.  A union bound excludes all fixed imbalances.  A labelled-slot
amplifier then preserves balance with high probability.
\end{program}

\section{Slowly growing support}

Let $S_m$ be a consecutive deficit support whose size tends to infinity, and
let $D_m(\delta)$ be its support loss.

\begin{conjecture}[Uniform truncation]
There exist constants $c,C>0$ such that
\[
 \sup_{\delta\in[0,1]}D_m(\delta)\le Ce^{-cm^2}.
\]
\end{conjecture}

This is suggested by the Gaussian tail structure of the selected full-support
deficit distribution.

\begin{conjecture}[Slow-support gap]
\label{conj:slow-support}
There exists an explicit sequence $m_n\to\infty$ and admissible signed profiles
supported on $m_n$ consecutive sizes such that
\[
 \chi(G_n)-\zeta(G_n)
 \ge
 \left(\frac{q^3}{4}-o(1)\right)H_n
\]
with high probability.
\end{conjecture}

The limiting coefficient is
\[
 \frac{q^3}{4}=0.083256162997232\ldots.
\]
Midpoint placement alone would give $q^3/8$.

\begin{program}[Dimension-uniform second moment]
Derive an explicit complexity envelope $K(m)$ for the endpoint and residual
estimates.  Prove the bare-skeleton and q-only errors in the form
\[
 \exp\left\{O\left(K(m)\frac n{N^4}\right)\right\}.
\]
Choose $m=m_n$ so that $K(m_n)=o(N)$.  If $K$ is polynomial, one may take
$m_n=\lfloor\log\log n\rfloor$; if $K$ is moderately exponential, a slower
sequence such as $\lfloor\log\log\log n\rfloor$ may be required.
\end{program}

\section{Matching upper bound and the full order conjecture}

\begin{conjecture}[Heckel order conjecture]
\[
 \chi(G_n)-\zeta(G_n)
 =\Theta\left(\frac n{(\log n)^3}\right)
\]
with high probability.
\end{conjecture}

The candidate proof addresses the lower-bound direction.  A matching upper
bound should be decomposed into a lower location theorem for $\zeta$ and a
third-order upper construction for $\chi$.

\begin{conjecture}[Cochromatic location corridor]
There is a full signed-profile root $r_{\mathrm{sgn}}(n)$ such that
\[
 r_{\mathrm{sgn}}(n)-o(H_n)
 \le\zeta(G_n)
 \le\rco(n)+o(H_n)
\]
with high probability.
\end{conjecture}

The lower inequality should follow from a first-moment union bound over all
signed profiles.  The upper inequality is the present constructive result.

\begin{program}[Full upper bound]
\begin{enumerate}[label=(\roman*)]
\item Determine the variational optimum over all signed profiles and prove a
uniform first-moment exclusion below it.
\item Construct ordinary colorings at $\rplus+O(H_n)$ rather than only at
first-order or $o(n/N^2)$ precision.
\item Compare the two phase-resolved roots to obtain
$\chi-\zeta=O(H_n)$.
\end{enumerate}
\end{program}

\section{General density and the two-layer model}

For $G(n,p)$, the one-class signed reward is
\[
 g_p(s)=p^{\binom s2}+(1-p)^{\binom s2}.
\]
At $p=1/2$, this reduces to $2^{1-\binom s2}$.  For $p\ne1/2$, the sign
variable has a size-dependent external field, and the overlap compatibility
sum becomes an inhomogeneous Ising-type partition function.

\begin{program}[Fixed $p\ne1/2$]
Determine the optimal clique proportion, phase-resolved signed root, and
whether a polynomial chromatic--cochromatic gap persists for every fixed
$p\in(0,1)$.
\end{program}

For the two-layer model, sample independent $G_1,G_2\sim G(n,1/2)$ and let
$X(G_1,G_2)$ be the minimum number of classes, each independent in at least one
layer.  For a fixed partition and layer assignment,
\[
 \Pp\{\text{all constraints hold}\}
 =2^{-\sum_i\binom{|C_i|}{2}},
\]
so its first moment is exactly the same as the signed one-graph witness after
summing the $2^k$ assignments.

\begin{program}[Two-layer comparison]
Obtain and verify the McDiarmid coupling reported in the Erd\H{o}s Problems
discussion, with explicit attribution.  Then calculate the two-layer second
moment and compare its constant and complexity with the one-graph cycle-space
formula.
\end{program}

\section{Recommended order of attack}

\begin{enumerate}[label=\arabic*.]
\item Close the endpoint physical equivalence and global deficit reindexing.
\item Promote Theorem~\ref{thm:phase-gap}, Corollary~\ref{cor:uniform}, and
Corollary~\ref{cor:complement} into the main paper.
\item Prove the near-root placement theorem, beginning with the explicit choice
$\theta_n=N^{-1/2}$.
\item Prove balance stability.
\item Certify the exact four-support phase minimum.
\item Choose between the slowly growing support program and the matching upper
bound as the main separate-paper project.
\end{enumerate}

\end{document}
